\chapter{Аналитическая часть}
\section{Расстояние Левенштейна}
\textbf{Расстояние Левенштайна} -- это метрика, используемая для определения минимального количества элементарных операций, необходимых для преобразования одной строки (последовательности символов) в другую. К таким операциям относятся: вставка, удаление и замена символа.

Расстояние Левенштейна широко применяется в задачах, связанных с обработкой текста и последовательностей данных. Например, в нечетком поиске, оно позволяет определить степень схожести строк, что помогает при поиске слов с опечатками. В коррекции орфографии расстояние Левенштейна используется для нахождения слов, которые максимально близки к заданному слову с ошибкой. В области обработки естественного языка (NLP) метрика помогает анализировать и сопоставлять предложения, тексты и другие текстовые структуры, чтобы оценить их схожесть.

Кроме того, расстояние Левенштейна находит применение в биоинформатике для сравнения последовательностей ДНК и белков, а также в системах контроля версий для отслеживания различий между файлами и их изменениями.\cite{lev}

\subsection{Рекуррентная формула расстояния Левенштейна}

Пусть существует две последовательности символов $S_1$ и $S_2$ длины N и M соответственно. Введем функцию $D(i, j)$, равную расстоянию Левенштейна между подстроками $S_1[1...i]$ и $S_2[1...j]$.

Тогда расстояние Левенштайна для двух строк может быть выражено следующей рекуррентной формулой:

\begin{equation}
	\label{eq:D}
	D(i, j) = 
	\begin{cases}
		0 &\text{i = 0 и j = 0}\\
		i &\text{i > 0 и j = 0}\\
		j &\text{i = 0 и j > 0}\\
		min \begin{cases}
			D(i, j - 1) + 1\\
			D(i - 1, j) + 1\\
			D(i - 1, j - 1) + f_e(S_1[i], S_2[j])\\
		\end{cases} &\text{i > 0 и j > 0}\\
	\end{cases}
\end{equation}
При этом
\begin{equation}
	\label{eq:f_e}
	f_e(a1, a2) = 
	\begin{cases}
		0 &\text{a1 == a2}\\
		1 &\text{a1 != a2}\\
	\end{cases}
\end{equation}

Из функции видно, что она вычисляет минимальное расстояние Левенштейна на текущем этапе (по индексу в первой строке), постепенно приводя первую строку к виду второй с конца. В рекурсивной части: первый случай символизирует вставку символа из второй строки в первую, второй — удаление символа из первой строки, а третий — замену символа первой строки на символ второй. Если символы уже совпадают, замена не требуется и её стоимость равна 0.\cite{algs}

\subsection{Алгоритм Левенштейна с кешированием}
Для избежания многократного вычисления $D(i, j)$ для одинаковых аргументов, используют итерационный алгоритм, который сохраняет промежуточные результаты в матрице размером $(N+1) \times (M+1)$.

Каждая ячейка матрицы $[i, j]$ хранит значение $D(i, j)$, представляющее расстояние Левенштейна между подстроками $S_1[1...i]$ и $S_2[1...j]$. Нулевая строка и нулевой столбец матрицы заполняются от 0 до соответствующего индекса, так как для преобразования пустой строки в строку длины $i$ требуется $i$ операций вставки.

Затем алгоритм построчно заполняет матрицу по формуле (\ref{eq:D}), используя уже сохранённые значения из матрицы, что исключает необходимость пересчёта значений $D$. 

Для оптимизации памяти на каждой итерации сохраняют только текущую строку и предыдущую, так как остальные строки не участвуют в дальнейшем вычислении.\cite{algs}

\section{Алгоритм Дамерау-Левенштейна}

\textbf{Расстояние Дамерау-Левенштайна} -- метрика, которая измеряет расстояние между двумя последовательностями символов, определяя минимальное количество операций (вставка, замена, удаление и транспозиция), необходимых для преобразования одной последовательности в другую. Транспозиция подразумевает перестановку двух соседних символов. В отличие от расстояния Левенштейна, эта метрика включает операцию транспозиции, что может значительно уменьшить итоговое расстояние между последовательностями.

\subsection{Рекурсивная формула расстояния Дамерау-Левенштейна}

Введем также как и для расстояния Левенштейна функцию $D(i, j)$, при этом с учётом транспозиции формула будет выглядеть так:

\begin{equation}
	\label{eq:Dd}
	D(i, j) = 
	\begin{cases}
		0 &\text{i = 0, j = 0}\\
		i &\text{i > 0, j = 0}\\
		j &\text{i = 0, j > 0}\\
		min \begin{cases}
			D(i, j - 1) + 1\\
			D(i - 1, j) + 1\\
			D(i - 1, j - 1) + f_e(S_1[i], S_2[j])\\
			D(i - 2, j - 2) + 1
		\end{cases} & \begin{aligned}
		&\text{i > 1,  j > 1,} \\ 
		&\text{$S_1$[i] = $S_2$[j - 1],} \\
		&\text{$S_1
  $[i - 1] = $S_2$[j]} \\
		\end{aligned} \\
		min \begin{cases}
			D(i, j - 1) + 1\\
			D(i - 1, j) + 1\\
			D(i - 1, j - 1) + f_e(S_1[i], S_2[j])\\
		\end{cases} & \text{Иначе}\\
	\end{cases}
\end{equation}

Фактическое отличие от алгоритма расстояния Левенштейна в том, что при j > 1 и i > 1, а также при условии что соседние символы в строках равны \guillemotleftнакрест\guillemotright, может быть проведена операция транспозиции, которая может сократить расстояние.

\subsection{Алгоритм Дамерау-Левенштейна с кешированием}
Существует итерационный алгоритм с кешированием, сохраняющий промежуточные значения в матрице либо только текущую и последние две строки. Алгоритм схож с алгоритмом для вычисления расстояния Левенштейна, за исключением того, что при совпадении двух соседних символов строк \guillemotleftнакрест\guillemotright\space(если $S_1[i] = S_2[j-1]$ и $S_1[i-1] = S_2[j]$) и при $i > 1$ и $j > 1$, необходимо учитывать вариант с транспозицией при расчёте текущей ячейки матрицы. Это позволяет учесть операцию перестановки символов, что может уменьшить итоговое расстояние.\cite{algs}

\section*{Вывод}
В результате аналитического раздела были рассмотрены расстояния Левенштейна и Дамерау-Левенштейна, а также их рекурсивные формулы, итерационные алгоритмы с кешированием матрицами и отдельными строками матриц.

\clearpage